\documentclass[a4paper,12pt,twoside]{report} % Dùng 12pt làm mặc định
\usepackage{scrextend}
\changefontsizes{13pt} % Ép lên 13pt bằng scrextend
\usepackage{mathptmx} % Đưa lên trước để set font chính xác
\usepackage[T5]{fontenc} % Bắt buộc T5 cho tiếng Việt thay vì T1
\usepackage[utf8]{inputenc}
\usepackage[vietnamese]{babel}
\usepackage[top=2cm, bottom=2cm, left=3.5cm, right=2.5cm]{geometry}
\usepackage{xurl}
\usepackage{appendix}
\usepackage{xcolor}
\usepackage{outlines}
\usepackage{graphicx}
\usepackage{fancybox}
\usepackage{indentfirst}
\usepackage{amsthm}
\usepackage{latexsym}
\usepackage{amsmath}
\usepackage{amssymb}
\usepackage{amsbsy}
\usepackage{array}
\usepackage{enumitem}
\usepackage{subfiles}
\usepackage{titlesec}
\usepackage{titletoc}
\usepackage{chngcntr}
\usepackage{pdflscape}
\usepackage{afterpage}
\usepackage[ruled,vlined]{algorithm2e}
\usepackage{capt-of}
\usepackage{multirow}
\usepackage{fancyhdr}

\usepackage[font=small,labelfont=bf]{caption}

\usepackage{listings}
\usepackage{float}
\usepackage{subcaption}

\usepackage[nonumberlist, nopostdot, nogroupskip, acronym]{glossaries}
\usepackage{glossary-superragged}
\setglossarystyle{superraggedheaderborder}
\usepackage{setspace}
\usepackage{parskip}

\usepackage{tocbasic}
\usepackage{blindtext}
\usepackage{csquotes}

\usepackage{tikz}
\usetikzlibrary{calc}


% ===================================================

\renewcommand{\bibname}{Danh_sach_tai_lieu_tham_khao} 
\usepackage[backend=biber,style=ieee]{biblatex}

\addbibresource{Danh_sach_tai_lieu_tham_khao.bib}

\input{lstlisting}

\newacronym{c2c}{C2C}{Customer-to-Customer}
\newacronym{tmdt}{TMĐT}{Thương mại điện tử}
\newacronym{hust}{HUST}{Hanoi University of Science and Technology}
\newacronym{escrow}{Escrow}{Thanh toán trung gian}
\newacronym{api}{API}{Application Programming Interface}
\newacronym{orm}{ORM}{Object-Relational Mapping}
\newacronym{sql}{SQL}{Structured Query Language}
\newacronym{jwt}{JWT}{JSON Web Token}
\newacronym{css}{CSS}{Cascading Style Sheets}
\newacronym{html}{HTML}{HyperText Markup Language}
\newacronym{jsx}{JSX}{JavaScript XML}
\newacronym{rsc}{RSC}{React Server Components}
\newacronym{ssr}{SSR}{Server-Side Rendering}
\newacronym{fcp}{FCP}{First Contentful Paint}
\newacronym{acid}{ACID}{Atomicity, Consistency, Isolation, Durability}
\newacronym{ipn}{IPN}{Instant Payment Notification}
\newacronym{url}{URL}{Uniform Resource Locator}
\newacronym{otp}{OTP}{One-Time Password}
\newacronym{https}{HTTPS}{Hypertext Transfer Protocol Secure}
\newacronym{ui}{UI}{User Interface}
\newacronym{ux}{UX}{User Experience}
\newacronym{oauth}{OAuth}{Open Authorization}
\newacronym{json}{JSON}{JavaScript Object Notation}
\newacronym{erd}{ERD}{Entity-Relationship Diagram}
\newacronym{dfd}{DFD}{Data Flow Diagram}
\newacronym{vnpay}{VNPAY}{Vietnam Payment Gateway}
\newacronym{ghtk}{GHTK}{Giao Hàng Tiết Kiệm}
\newacronym{ghn}{GHN}{Giao Hàng Nhanh}
\newacronym{oop}{OOP}{Object-Oriented Programming}

\makenoidxglossaries


% ===================================================


\fancypagestyle{plain}{%
\fancyhf{}
\fancyfoot[RO,RE]{\thepage}
\renewcommand{\headrulewidth}{0pt}
\renewcommand{\footrulewidth}{0pt}}

\setlength{\headheight}{10pt}

\def \TITLE{ĐỒ ÁN TỐT NGHIỆP}
\def \PROJECTTITLE{PHÁT TRIỂN NỀN TẢNG MUA BÁN SÁCH C2C HỖ TRỢ THANH TOÁN TRUNG GIAN}
\def \AUTHOR{Nguyễn Xuân Hoàng}
\def \STUDENTID{20210149P}
\def \INSTRUCTOR{TS. Nguyễn Tiến Thành}

% ===================================================
\titleformat{\chapter}[hang]{\centering\bfseries}{CHƯƠNG \thechapter.\ }{0pt}{}[]

\titleformat 
    {\chapter} % command
    [hang] % shape
    {\centering\bfseries} % format
    {CHƯƠNG \thechapter.\ } % label
    {0pt} %sep
    {} % before
    [] % after
\titlespacing*{\chapter}{0pt}{-20pt}{20pt}

\titleformat
    {\section} % command
    [hang] % shape
    {\bfseries} % format
    {\thechapter.\arabic{section}\ \ \ \ } % label
    {0pt} %sep
    {} % before
    [] % after
\titlespacing{\section}{0pt}{\parskip}{0.5\parskip}

\titleformat
    {\subsection} % command
    [hang] % shape
    {\bfseries} % format
    {\thechapter.\arabic{section}.\arabic{subsection}\ \ \ \ } % label
    {0pt} %sep
    {} % before
    [] % after
\titlespacing{\subsection}{30pt}{\parskip}{0.5\parskip}

\renewcommand\thesubsubsection{\alph{subsubsection}}
\titleformat
    {\subsubsection} % command
    [hang] % shape
    {\bfseries} % format
    {\alph{subsubsection}, \ } % label
    {0pt} %sep
    {} % before
    [] % after
\titlespacing{\subsubsection}{50pt}{\parskip}{0.5\parskip}

% ===================================================
\usepackage{hyperref}
\hypersetup{pdfborder = {0 0 0}}
\hypersetup{pdftitle={\TITLE},
	pdfauthor={\AUTHOR}}
	
\usepackage[all]{hypcap}

\graphicspath{{figures/}{../figures/}}

\counterwithin{figure}{chapter}

\title{\bf \TITLE}
\author{\AUTHOR}

\setcounter{secnumdepth}{3}

\theoremstyle{definition}
\newtheorem{example}{Ví dụ}[chapter]

\onehalfspacing
\setlength{\parskip}{6pt}
\setlength{\parindent}{15pt}



% =========================== BODY ===============
\begin{document}
\subfile{Bia}
\subfile{Bia_lot}

% ===================================================
\pagestyle{empty}
\newpage
\pagenumbering{gobble}
\subfile{Chuong/0_Phieu_giao_nhiem_vu}
\clearpage
\subfile{Chuong/0_Nhan_xet}
\clearpage
\subfile{Chuong/0_Loi_cam_doan}
\clearpage
\subfile{Chuong/0_Loi_cam_on}
\clearpage
\subfile{Chuong/0_Tom_tat}
\clearpage
\subfile{Chuong/0_Abstract_EN}

% ===================================================
\newpage
\pagenumbering{gobble}
\renewcommand*\contentsname{MỤC LỤC}
\titlecontents{chapter}
    [0.0cm]             % left margin
    {\bfseries\vspace{0.3cm}}                  % above code
    {{\bfseries{\scshape} CHƯƠNG \thecontentslabel.\ }} % numbered format
    {}         % unnumbered format
    {\titlerule*[0.3pc]{.}\contentspage}
    
\titlecontents{section}
    [0.0cm]             % left margin
    {\vspace{0.3cm}}                  % above code
    {\thecontentslabel \ } % numbered format
    {}         % unnumbered format
    {\titlerule*[0.3pc]{.}\contentspage}
    
\titlecontents{subsection}
    [1.0cm]             % left margin
    {\vspace{0.3cm}}                  % above code
    {\thecontentslabel \ } % numbered format
    {}         % unnumbered format
    {\titlerule*[0.3pc]{.}\contentspage}

\addtocontents{toc}{\protect\thispagestyle{empty}}
\tableofcontents 
\thispagestyle{empty}
\clearpage

\pagenumbering{roman}
\renewcommand{\listfigurename}{DANH MỤC HÌNH VẼ}
{\let\oldnumberline\numberline
\renewcommand{\numberline}{Hình~\oldnumberline}
\listoffigures} 
\newpage


\renewcommand{\listtablename}{DANH MỤC BẢNG BIỂU}
{\let\oldnumberline\numberline
\renewcommand{\numberline}{Bảng~\oldnumberline}
\listoftables}

\glsaddall 
\renewcommand*{\acronymname}{DANH MỤC THUẬT NGỮ VÀ TỪ VIẾT TẮT}
\renewcommand*{\entryname}{Thuật ngữ}
\renewcommand*{\descriptionname}{Ý nghĩa}
\printnoidxglossary[type=\acronymtype]

\renewcommand\appendixname{PHỤ LỤC}
\renewcommand\appendixpagename{PHỤ LỤC}
\renewcommand\appendixtocname{PHỤ LỤC}



% ===================================================


\newpage
\pagenumbering{arabic}

\pagestyle{fancy}
\fancyhf{}
\fancyhead[RE, LO]{\leftmark}
\fancyfoot[RE, LO]{\thepage}

\chapter{GIỚI THIỆU CHUNG}
\subfile{Chuong/1_Gioi_thieu}

\newpage
\chapter{CƠ SỞ LÝ THUYẾT VÀ CÔNG NGHỆ SỬ DỤNG}
\subfile{Chuong/2_Co_so_ly_thuyet}

\newpage
\chapter{PHÂN TÍCH YÊU CẦU VÀ GIẢI PHÁP ĐỀ XUẤT}
\subfile{Chuong/3_Giai_phap_de_xuat}

\newpage
\chapter{THIẾT KẾ HỆ THỐNG}
\subfile{Chuong/4_Thiet_ke_he_thong}

\newpage
\chapter{TRIỂN KHAI VÀ ĐÁNH GIÁ KẾT QUẢ}
\subfile{Chuong/5_Thuc_nghiem_ket_qua}

\newpage
\chapter{KẾT LUẬN VÀ HƯỚNG PHÁT TRIỂN}
\subfile{Chuong/6_Ket_luan}

\newpage
\chapter*{MỘT SỐ LƯU Ý VỀ TÀI LIỆU THAM KHẢO}
\label{chapter:reference}
\subfile{Chuong/7_Ghi_chu_tai_lieu}

% ===================================================
\newpage
\renewcommand\bibname{TÀI LIỆU THAM KHẢO}
\printbibliography
\phantomsection\addcontentsline{toc}{chapter}{TÀI LIỆU THAM KHẢO}

\appendixpage
\appendix
\addappheadtotoc

\titleformat{\chapter}[hang]{\centering\bfseries}{ \thechapter.\ }{0pt}{}[]
\titlespacing*{\chapter}{0pt}{-20pt}{20pt}

\titlecontents{chapter}
    [0.0cm]             % left margin
    {\bfseries\vspace{0.3cm}}                  % above code
    {{\bfseries{\scshape} \thecontentslabel.\ }} % numbered format
    {}         % unnumbered format
    {\titlerule*[0.3pc]{.}\contentspage} 
    
\chapter{HƯỚNG DẪN CÀI ĐẶT VÀ SỬ DỤNG}
\subfile{Chuong/Phu_luc_A}

\newpage
\chapter{ĐẶC TẢ USE CASE CHI TIẾT}
\subfile{Chuong/Phu_luc_B}

\newpage
\chapter{HÌNH ẢNH GIAO DIỆN HỆ THỐNG}
\subfile{Chuong/Phu_luc_C}

\end{document}